\chapter{Future Work}
\lhead{Chapter 8. \emph{Future Work}}
\label{chap.8}

While Q-SFTP provides a comprehensive post-quantum file transfer solution, several promising directions for future enhancement have been identified:

\begin{itemize}
    \item \textbf{Hybrid Key Exchange:} Combining Kyber512 with X25519 (classical ECDH) in a hybrid KEM would provide defense-in-depth against potential unforeseen weaknesses in lattice-based schemes. This approach aligns with IETF draft recommendations for hybrid post-quantum TLS \cite{chen2022hybridtls}.
    
    \item \textbf{End-to-End File Encryption:} Extending the system with per-file encryption keys derived from user passwords or client-side key pairs, ensuring that even the server operator cannot access file contents (zero-knowledge storage).
    
    \item \textbf{Formal Security Verification:} Applying symbolic model checking tools such as ProVerif or the Tamarin prover to formally verify the Q-TLS handshake protocol's security properties (secrecy, authentication, forward secrecy) against automated adversary models.
    
    \item \textbf{WebSocket Transport:} Migrating the chunked transfer protocol from HTTP polling to WebSocket for reduced latency, bidirectional communication, and server-push capabilities — particularly beneficial for real-time transfer progress updates and large-file downloads.
    
    \item \textbf{Performance Benchmarking:} Systematic evaluation of handshake latency, file transfer throughput, and memory usage across varying network conditions (latency, bandwidth, packet loss), file sizes, and concurrent user loads to establish quantitative performance baselines.
    
    \item \textbf{Multi-Server and Cloud Deployment:} Support for distributed deployment with load balancing, database replication, and shared-nothing architecture for horizontal scalability in enterprise environments.
    
    \item \textbf{Advanced Integrity Features:} Extending the hash registry with version tracking (file history), Merkle tree-based bulk verification for efficient integrity checking of large repositories, and real-time file change notifications.
    
    \item \textbf{Cross-Platform Client:} Developing native desktop clients (Electron or Qt-based) with drag-and-drop support and system tray integration, providing an experience comparable to classical FTP clients like FileZilla.
    
    \item \textbf{Post-Quantum Algorithm Agility:} Implementing a negotiation mechanism in the handshake that allows the client and server to dynamically select from multiple PQC algorithms (e.g., Kyber, NTRU, Classic McEliece) based on security requirements and resource constraints.
    
    \item \textbf{Quantum Key Distribution (QKD) Integration:} As quantum networks mature, exploring the integration of hardware-based QKD for session key establishment as an alternative or supplement to software-based PQC key encapsulation.
\end{itemize}

These enhancements will further strengthen Q-SFTP's position as a practical, scalable, and forward-compatible standard for secure file transfer in the post-quantum era.
